\section{Technical Foundations and Infrastructure}\label{sec:implementation-foundations}

The implementation of the DSL compiler is a critical engineering effort that balances high-level musical
expressiveness with low-level binary precision.
This section details the selection of the technology stack and the structural organization of the project,
providing a rationale for the architectural decisions that support the thesis's goals.

\subsection{Language Choice: Modern C++ (C++23)}

The compiler is written in modern C++, specifically targeting the C++23 standard.
While higher-level languages such as Python or JavaScript are commonly used for DSL development, C++ was
selected for three primary reasons:
\begin{enumerate}
  \item \textbf{Deterministic Performance:} Musical temporal resolution requires high precision. By using
    C++, the compiler can perform recursive unrolling and temporal math with minimal overhead, achieving the
    "insane" compilation speeds necessary for real-time feedback in the playground.
  \item \textbf{Robust Type System:} Modern C++ provides a wealth of tools for building type-safe compilers.
    Features such as \texttt{std::variant} and \texttt{std::optional} allow for a functional programming
    style that is ideally suited for processing tree-like structures (ASTs).
  \item \textbf{Binary Mastery:} Since the final build target is a low-level MIDI binary file, C++'s native
    support for bitwise operations and big-endian serialization makes it the most natural choice for
    implementing the backend.
\end{enumerate}

\subsubsection{Key C++23 Features in the Architecture}

The codebase leverages several advanced features to simplify compiler logic:
\begin{itemize}
  \item \textbf{\texttt{std::variant} as a Sum Type:} Every AST node and IR value is represented as a
    variant. This allows the compiler to treat different musical constructs (e.g., a Note vs. a Loop) as a
    single type while enforcing exhaustive handling during traversal via \texttt{std::visit}.
  \item \textbf{Value-Oriented Memory Management:} By combining \texttt{std::variant} with
    \texttt{std::unique\_ptr}, the compiler avoids the "Object Soup" anti-pattern. Nodes are value-like, and
    recursion is handled via smart pointers that guarantee zero memory leaks without the non-deterministic
    latency of a garbage collector.
  \item \textbf{Standard Ranges and Views:} Used extensively in the backend and lowering passes to sort
    musical events and filter diagnostic logs efficiently.
\end{itemize}

\subsection{Modular Build Infrastructure (CMake)}

To ensure scalability and maintainability, the project is structured as a modular hierarchy managed by CMake.
The source code is partitioned into discrete components, each with its own responsibilities and boundaries.

\subsubsection{Library Decomposition}

The compiler logic is distributed across five static libraries:
\begin{enumerate}
  \item \textbf{\texttt{dsl\_common}:} The foundational library containing the core musical enums
    (\texttt{Pitch}, \texttt{Accidental}), the \texttt{Note} structure, and the \texttt{DiagnosticsEngine}.
    It is linked by every other module.
  \item \textbf{\texttt{dsl\_frontend}:} Wraps the Flex-generated scanner and Bison-generated parser. It
    depends on \texttt{dsl\_common} for location tracking and note literals.
  \item \textbf{\texttt{dsl\_semantic}:} Contains the \texttt{Traversal} engine and symbol table logic. It
    takes a raw AST and produces an \texttt{AnalysisResult}.
  \item \textbf{\texttt{dsl\_lowerer}:} The unrolling engine. It transforms an \texttt{AnalysisResult} into a
    flat list of \texttt{NoteEvent} objects.
  \item \textbf{\texttt{dsl\_backend}:} The MIDI serialization module. It handles the binary translation of
    the linear event stream.
\end{enumerate}

\subsubsection{Directory Organization and Public API}

The project follows the "Pitch-Perfect" directory convention:
\begin{itemize}
  \item \textbf{\texttt{include/dsl/}:} Contains only the public headers (e.g., \texttt{compiler.hpp}). This
    defines the interface for external consumers, such as the CLI wrapper or the playground backend.
  \item \textbf{\texttt{src/dsl/}:} Contains the private implementation details, including the recursive
    visitors and generated Flex/Bison sources.
\end{itemize}
This separation ensures that the internal complexity of the compiler pipeline is hidden behind a simple, high-level API.
The top-level \texttt{dsl::compile} function serves as the entry point, orchestrating the sequential flow of
characters to tree, tree to timeline, and timeline to MIDI.
